\documentclass[../DoAn.tex]{subfiles}
\begin{document}

Chương~4 đã trình bày tổng quan thiết kế kiến trúc và kết quả thực nghiệm. Chương này
đi sâu vào năm đóng góp kỹ thuật cốt lõi của đồ án, mỗi mục trình bày rõ bài toán,
giải pháp và kết quả đạt được.

% ============================================================
\section{Hệ thống đọc bản đồ động và dựng lưới thống nhất}
\label{section:5.1}

\subsection*{Bài toán}
Game thường bake cứng bản đồ vào scene Unity, khiến việc thay đổi bản đồ đòi hỏi
mở lại scene và chỉnh sửa thủ công. Khi cần chạy backtest trên 5 bản đồ khác nhau
với kịch bản MAPF, cách tiếp cận này không khả thi.

\subsection*{Giải pháp}
Lớp \texttt{MapLoader} đọc tệp \texttt{.map} theo chuẩn MovingAI tại thời điểm
chạy (runtime): phân tích header (\texttt{height}, \texttt{width}) và ma trận ký tự
(\texttt{.} = ô đi được, \texttt{@} = tường), sau đó tạo \texttt{GameObject} cho mỗi
ô theo đúng tọa độ. Ô tường nhận \texttt{BoxCollider2D} trên layer~\texttt{ObstaclesMovement}
(chặn vật lý) và trigger collider trên layer~\texttt{Hittable} (nhận đạn). Bốn tường
biên được tạo tự động bao quanh bản đồ.

Hệ thống expose ba API thống nhất mà mọi thuật toán MAPF và spawner đều dùng:
\texttt{IsWalkable(cell)}, \texttt{CellToWorld(cell)}, \texttt{WorldToCell(pos)}.
Tọa độ ô dùng quy ước góc trên-trái, Y tăng xuống dưới (khớp với hàng trong tệp
\texttt{.map}); tọa độ thế giới Unity tâm ở $(0,0)$, trục XY.

\subsection*{Kết quả}
Chuyển đổi bản đồ trong backtest chỉ tốn vài mili-giây (không cần reload scene).
Cùng một bản đồ \texttt{.map} chạy được trên cả ba thuật toán mà không cần chỉnh sửa.
Hệ thống đã xử lý thành công 5 bản đồ kích thước từ 32×32 đến 251×180 ô, bao gồm
cả bản đồ mê cung 128×128 với 14\,818 ô đi được.

% ============================================================
\section{Triển khai A* thời gian thực với tái hoạch định định kỳ}
\label{section:5.2}

\subsection*{Bài toán}
A* truyền thống tính một đường đi tĩnh từ điểm đầu đến đích. Trong game thời gian thực,
đích có thể thay đổi (Eagle di chuyển hoặc bị che khuất), bản đồ có thể thay đổi (thùng
cản mới xuất hiện), và agent bị cản bởi agent khác. Cần cơ chế tái hoạch định tự động
mà không làm giật game.

\subsection*{Giải pháp}
\texttt{GridAStarPathfinder} thực thi A* 4 láng giềng với heuristic Manhattan trên
ma trận \texttt{char[][]}, tái sử dụng \texttt{List<Vector2Int>} để tránh garbage
collection. \texttt{GridEnemyAgent} quản lý vòng lặp điều hướng: tái hoạch định sau
mỗi \texttt{replanInterval} giây (mặc định), chuyển sang chế độ bắn khi Eagle trong
tầm và line-of-sight, phục hồi khi bị kẹt quá 2 giây.

Ngưỡng tiếp cận waypoint dùng dot-product 0.96 (thay vì khoảng cách Euclidean) để
tránh rung lắc ở các waypoint gần: agent chỉ chuyển sang waypoint tiếp theo khi hướng
di chuyển gần như song song với hướng đến waypoint. Hình~\ref{fig:seq_astar} mô tả luồng
tương tác giữa agent, bộ tìm đường và bộ điều khiển trong một chu kỳ tái hoạch định.

\begin{figure}[htbp]
    \centering
    \includegraphics[width=0.95\textwidth]{Hinhve/seq_astar_replan.png}
    \caption{Biểu đồ trình tự chu kỳ tái hoạch định của A*}
    \label{fig:seq_astar}
\end{figure}

\subsection*{Kết quả}
Thực nghiệm trên Alpha32 (90\% ô đi được, 6 agent): thời gian hoàn thành trung bình
35,1 giây, 159 lần replan. Thuật toán chạy ổn định 6 lần lặp với độ lệch nhỏ
(< 0,5 giây), xác nhận tính tất định của A* trên cùng bản đồ.

% ============================================================
\section{Tích hợp PIBT phối hợp đa agent không va chạm}
\label{section:5.3}

\subsection*{Bài toán}
Khi nhiều agent A* chạy độc lập, chúng thường cố tranh vào cùng một ô tại cùng thời
điểm, gây va chạm ở lớp vật lý và replan phản ứng liên tục. Cần cơ chế phối hợp ở
lớp lập kế hoạch để loại bỏ xung đột từ gốc.

\subsection*{Giải pháp}
\texttt{MapScenarioBootstrapPIBT} tạo một bộ giải PIBT dùng chung cho tất cả agent.
Tại mỗi tick lập kế hoạch, bộ giải nhận vị trí hiện tại của tất cả agent, chạy một
bước PIBT, và trả về ô di chuyển tiếp theo cho từng agent. Agent nhận lệnh di chuyển
qua \texttt{GridEnemyAgentPIBT.SetNextWaypoint()} và thực thi qua \texttt{TankController}
-- cùng pipeline với A*.

Cơ chế kế thừa ưu tiên đảm bảo nếu agent ưu tiên thấp bị agent ưu tiên cao chặn ô,
nó tạm nhường ô đó và chờ hoặc tìm ô thay thế. Nếu không có ô nào khả dụng (deadlock
cục bộ), thuật toán hoàn tác (backtracking) lên agent ưu tiên cao để thử phương án khác.
Hình~\ref{fig:seq_pibt} mô tả luồng một tick lập kế hoạch của bộ giải PIBT dùng chung.

\begin{figure}[htbp]
    \centering
    \includegraphics[width=0.95\textwidth]{Hinhve/seq_pibt_csharp.png}
    \caption{Biểu đồ trình tự một tick lập kế hoạch của PIBT C\#}
    \label{fig:seq_pibt}
\end{figure}

\subsection*{Kết quả}
PIBT C\# hoàn thành nhiệm vụ trên 4/5 bản đồ với thời gian tương đương A* (chênh
lệch < 10\%). Số lần replan cao hơn A* 5--8 lần (do PIBT đếm mỗi bước phối hợp
là một replan), nhưng số lần phục hồi sau kẹt thấp hơn, cho thấy PIBT giảm được
tình trạng deadlock ở lớp vật lý.

% ============================================================
\section{Cầu nối TCP đến máy chủ PIBT C++}
\label{section:5.4}

\subsection*{Bài toán}
Thuật toán PIBT chuẩn được triển khai tối ưu bằng C++. Việc port toàn bộ sang C\#
sẽ mất nhiều công sức và dễ mất đi tối ưu hóa. Cần cách tích hợp thuật toán C++ vào
Unity mà không sửa engine game.

\subsection*{Giải pháp}
Giao thức TCP nội bộ tại localhost cổng 7777 sử dụng ba loại thông điệp JSON. Thông điệp
Hello được Unity gửi kèm kích thước bản đồ và số agent để server khởi tạo bộ lập kế hoạch
và xác nhận. Thông điệp PlanStep được Unity gửi kèm vị trí tất cả agent ở mỗi bước, và
server chạy một bước PIBT rồi trả về ô mục tiêu cho từng agent. Thông điệp Shutdown kết
thúc phiên khi ván chơi kết thúc. Phía Unity, lớp \texttt{PIBTTcpClient} quản lý kết nối
bất đồng bộ, timeout và tự động kết nối lại, còn \texttt{GridEnemyAgentPIBT\_TCP} nhận
lệnh từ client và điều khiển \texttt{TankController} như hai biến thể kia.
Hình~\ref{fig:seq_pibt_tcp} mô tả luồng trao đổi thông điệp giữa Unity và server C++.

\begin{figure}[htbp]
    \centering
    \includegraphics[width=0.95\textwidth]{Hinhve/seq_pibt_tcp.png}
    \caption{Biểu đồ trình tự trao đổi thông điệp PIBT qua TCP}
    \label{fig:seq_pibt_tcp}
\end{figure}

\subsection*{Kết quả}
PIBT C++/TCP hoàn thành nhiệm vụ trên cùng 4/5 bản đồ với thời gian tương đương
(chênh lệch < 5\% so với PIBT C\#). Agent TCP di qua nhiều ô hơn 20--70\% do độ
trễ TCP (một tick) khiến agent đôi khi chờ lệnh; đây là chi phí của kiến trúc
client--server mà có thể giảm bằng cách tăng tần suất tick. Kết nối TCP ổn định
trong suốt 90 run thực nghiệm mà không gặp lỗi mất kết nối.

% ============================================================
\section{Hệ thống backtest tự động với chướng ngại động}
\label{section:5.5}

\subsection*{Bài toán}
Để so sánh thuật toán một cách khách quan, cần chạy hàng chục kịch bản lặp lại trên
nhiều bản đồ và thuật toán, thu thập chỉ số nhất quán, và kiểm tra khả năng của hệ
thống khi môi trường thay đổi trong khi agent đang di chuyển. Chạy thủ công từng kịch
bản là không thực tế và không tái lập được.

\subsection*{Giải pháp}
\texttt{BacktestRunner} tự động lặp qua tất cả tổ hợp (map × algorithm × rep) theo
lịch trình định sẵn. Mỗi run: reset scene, spawn agent, chạy kịch bản, thu thập chỉ
số ở cấp run (summary) và cấp agent (per-agent CSV). Sau khi hết run, gọi script Python
để sinh biểu đồ HTML/PNG từ CSV.

\texttt{DynamicObstacleSpawner} bổ sung chế độ stress-test: thùng kim loại tối màu
liên tục spawn/despawn trực tiếp trên 1--6 ô phía trước đường đi của agent (ưu tiên
chặn ngay ô tiếp theo), ép agent phải tái hoạch định ngay lập tức.
\texttt{MapLoader.MarkCellBlocked/UnmarkCellBlocked} cập nhật ma trận grid đồng thời
để cả ba thuật toán đều nhận được cập nhật nhất quán.

\subsection*{Kết quả}
Hệ thống đã chạy thành công 90 run (môi trường tĩnh) + 90 run (môi trường động),
tổng cộng 180 run liên tục mà không crash. Khi bật chướng ngại động: số lần replan
tăng trung bình 50\% với A* và 30\% với PIBT, thời gian hoàn thành tăng $\approx$15\%.
Kết quả này xác nhận cả ba thuật toán đều phản ứng được với thay đổi môi trường
thời gian thực, và PIBT (nhờ cơ chế phối hợp) tái hoạch định ít hơn A* khi môi
trường bị xáo trộn.

\textbf{Kết chương~5:} Năm đóng góp kỹ thuật của đồ án -- từ hạ tầng bản đồ động
đến ba cấp độ thuật toán MAPF và hệ thống backtest -- tạo thành một nền tảng hoàn
chỉnh để nghiên cứu và so sánh thuật toán MAPF trong môi trường game thực tế. Kết
quả thực nghiệm được trình bày chi tiết trong Chương~4.

\end{document}
